\documentclass[11pt,a4paper]{article}
\usepackage[T1]{fontenc}
\usepackage[utf8]{inputenc}
\usepackage{lmodern}
\usepackage[margin=27mm,headheight=15pt]{geometry}
\usepackage{amsmath,amssymb,amsthm,mathtools}
\usepackage{booktabs,longtable,array,enumitem,microtype}
\usepackage{xcolor,fancyhdr}
\usepackage[unicode,colorlinks=true,linkcolor=blue!45!black,citecolor=blue!45!black,urlcolor=blue!45!black]{hyperref}
\usepackage{bookmark}
\definecolor{ink}{RGB}{25,42,64}
\definecolor{accent}{RGB}{35,91,120}
\newcommand{\R}{\mathbb R}
\newcommand{\N}{\mathbb N}
\newcommand{\E}{\mathbb E}
\newcommand{\Pbb}{\mathbb P}
\newcommand{\norm}[1]{\left\lVert#1\right\rVert}
\newcommand{\abs}[1]{\left\lvert#1\right\rvert}
\newcommand{\dd}{\,\mathrm d}
\DeclareMathOperator{\dist}{dist}
\DeclareMathOperator{\supp}{supp}
\DeclareMathOperator{\spanop}{span}
\DeclareMathOperator{\diam}{diam}
\DeclareMathOperator{\Lip}{Lip}
\newtheorem{theorem}{Theorem}[section]
\newtheorem{lemma}[theorem]{Lemma}
\newtheorem{proposition}[theorem]{Proposition}
\newtheorem{corollary}[theorem]{Corollary}
\theoremstyle{definition}
\newtheorem{definition}[theorem]{Definition}
\theoremstyle{remark}
\newtheorem{remark}[theorem]{Remark}
\numberwithin{equation}{section}
\setlist{topsep=4pt,itemsep=2pt,parsep=0pt}
\setlength{\emergencystretch}{2em}
\pagestyle{fancy}
\fancyhf{}
\fancyhead[L]{\small\color{ink}AIM 569}
\fancyhead[R]{\small 4 November 2026}
\fancyfoot[C]{\thepage}
\renewcommand{\headrulewidth}{0.3pt}


\urlstyle{same}
\title{Discontinuous information can improve randomized recovery on an incomplete normed space}
\author{Matthew J. Colbrook\thanks{The solutions were found by AI. The submitted arguments and the repository audit are AI work; independent human review and formal verification are not claimed.}\\[4pt]
\small Department of Applied Mathematics and Theoretical Physics\\
\small University of Cambridge, Cambridge CB3 0WA, United Kingdom\\
\small\href{mailto:m.colbrook@damtp.cam.ac.uk}{m.colbrook@damtp.cam.ac.uk}}
\date{4 November 2026}
\hypersetup{pdftitle={Discontinuous information can improve randomized recovery on an incomplete normed space},pdfauthor={Matthew J. Colbrook},pdfsubject={AIM 569: AI-found solution or partial result}}
\begin{document}
\maketitle
\begin{abstract}
On $F=(\ell_1,\|\cdot\|_2)$, for $S(x)=\|x\|_2^2$ on the unit ball, every finite expected-budget algorithm using continuous linear information has optimal root-mean-square error $1/2$. One discontinuous algebraic linear query attains error at most $\sqrt6-2<1/2$. This is a counterexample to AIM 569 as stated for normed spaces; no conclusion for Banach spaces alone is claimed.
\end{abstract}
\noindent\textbf{Catalogue target.} \href{https://github.com/MColbrook/AIM/blob/c4036871b0bf15cbc83e79c780b7191bcbb2cfed/problems/569-does-discontinuous-information-help-randomized-approximation.md}{AIM 569 at repository revision \texttt{c4036871b0bf}}. The package preserves its full statement in \texttt{statement.md}; the review and target comparison are linked from the package README.
\medskip

\section{Problem 569: a counterexample with one discontinuous query}
\label{sec:569}

The statement permits incomplete normed input spaces and continuous
nonlinear target maps. Both points matter here. Take
\[
 F=(\ell_1,\norm{\cdot}_2),\qquad G=\R,\qquad
 S(x)=\norm x_2^2,\qquad B=\{x\in\ell_1:\norm x_2<1\}.
 \tag{569.1}
\]
Thus the underlying vector space is \(\ell_1\), but its topology is
the one inherited from \(\ell_2\). Let $\Lambda_c$ be its continuous linear functionals and $\Lambda_a$ all algebraic linear functionals. In either class an algorithm may choose its next query and stopping rule from the preceding answers and a random seed. Its finite query count $N(x,\omega)$ and output $A_\omega(x)$ must be jointly measurable for the norm Borel structures and the seed sigma algebra. For $\Lambda\in\{\Lambda_c,\Lambda_a\}$ define
\[
e_n^{\rm ran}(S,\Lambda)=\inf_A\sup_{x\in B}\bigl(\E_\omega|S(x)-A_\omega(x)|^2\bigr)^{1/2},
\]
where $\sup_{x\in B}\E_\omega N(x,\omega)\le n$. No measurability of the history-to-query rule is added.

\begin{theorem}
\label{thm:569}
For (569.1), in the precise expected-budget RMS model of the problem,
\[
 e_n^{\rm ran}(S,\Lambda_c)=\frac12\qquad(n\ge1).
 \tag{569.2}
\]
On the other hand, with \(\kappa_n=\sqrt{2/(n+2)}\),
\[
 e_n^{\rm ran}(S,\Lambda_a)
           \le\frac{\kappa_n}{1+\kappa_n}\longrightarrow0 .
 \tag{569.3}
\]
In particular,
\[
 e_1^{\rm ran}(S,\Lambda_a)
          \le\sqrt6-2<\frac12
          =e_1^{\rm ran}(S,\Lambda_c).
 \tag{569.4}
\]
Consequently the universal equality in Problem 569 is false.
\end{theorem}

This settles the stated model. It does not assert a counterexample to
a version requiring the input space to be Banach. The distinction
between continuous and algebraic information discussed in
\cite{hnw2011} is retained exactly.

\subsection{An admissible algebraic estimator}

Let \(\Omega=\{-1,1\}^{n\times\N}\) carry product fair-sign measure, and
write \(r_{ij}\) for its coordinate signs. Define
\[
 L_i^\omega(x)=\sum_{j\ge1}r_{ij}(\omega)x_j,\qquad
 T_n(x,\omega)=\frac1n\sum_{i=1}^nL_i^\omega(x)^2 .
 \tag{569.5}
\]
For every \(x\in F\) and every seed, these series converge absolutely.
Thus each probe is an everywhere-defined algebraic real linear
functional. It is discontinuous for the chosen norm: a continuous
functional on the dense subspace \((\ell_1,\norm{\cdot}_2)\) extends
to \(\ell_2\), whereas its coefficient sequence
\((r_{ij})_j\) is not in \(\ell_2\).

Each probe is nevertheless jointly Borel in \(x\) and \(\omega\):
the finite partial sums are jointly continuous, and their pointwise
limit exists everywhere. Hence the estimator and its constant
cardinality \(n\) satisfy all stated measurability requirements.
There is no exceptional-input modification or algebraic-basis
extension.

For \(s=\norm x_2^2\), the independent rows satisfy
\[
 \E L_i^2=s,\qquad
 \E L_i^4=3s^2-2\sum_jx_j^4\le3s^2,\qquad
 \operatorname{Var}(T_n)\le \frac{2s^2}{n}.
 \tag{569.6}
\]
The fourth-moment identity follows first for finite sums and then by
\(L^4\) convergence of Rademacher series. In particular, the unbiased
estimator already gives error at most \(\sqrt{2/n}\).

For the stronger bound, set
\[
 a_n=\frac{\kappa_n}{1+\kappa_n},\qquad
 b_n=1-\kappa_n,\qquad A_n=a_n+b_nT_n .
 \tag{569.7}
\]
For arbitrary real \(a,b\), (569.6) gives
\[
 \E\abs{a+bT_n-s}^2
       \le \bigl(a+(b-1)s\bigr)^2+\frac{2b^2}{n}s^2 .
 \tag{569.8}
\]
This quadratic is convex in \(s\in[0,1]\), so its maximum is attained
at an endpoint. For a fixed \(a\), the endpoint at \(s=1\) is minimized
by \(b=n(1-a)/(n+2)\), with value
\(\kappa_n^2(1-a)^2\). Balancing this with the endpoint \(a^2\)
gives exactly (569.7), and both endpoint values equal \(a_n^2\).
This proves (569.3). For \(n=1\), \(a_1=\sqrt6-2\).
These constants optimize this affine-estimator calculation; no
optimality claim is made for all algebraic-information algorithms.

\subsection{A local weak-shift lemma}

The continuous-information lower bound requires a different
consequence of Gaussian conditioning.

\begin{lemma}[Insensitivity to weakly null shifts]
\label{lem:weakshift}
Let \(H\) be a real separable Hilbert space, and let \(F\subset H\)
be a dense Borel linear subspace with the inherited norm. Let \(\mu\)
be a centered Gaussian law on \(H\), with injective covariance \(C\),
such that
\[
                       \mu(F)=1,\qquad C(H)\subset F.
 \tag{W1}
\]
Let \(U\subset H\) be open. Suppose
\(A:U\cap F\to[-M,M]\) is Borel and uses at most \(k\) adaptive
continuous affine scalar measurements on \(F\).

Let \(w:H\to[0,\infty)\) be bounded and Borel, with
\[
             w(x)=0\quad\text{unless }
             \operatorname{dist}(x,H\setminus U)\ge R
 \tag{W2}
\]
for some \(R>0\). If \(v_j\in F\) is weakly null in \(H\) and
\(\limsup_j\norm{v_j}_H<R\), then, extending \(A\) by zero outside
\(U\cap F\),
\[
             \int_H w(x)\bigl[A(x+v_j)-A(x)\bigr]\,d\mu(x)
                        \longrightarrow0.
 \tag{W3}
\]
Again, no measurability of the adaptive decision rules is assumed.
\end{lemma}

\begin{proof}
The zero extension is Borel because \(F\) is Borel in \(H\).
For large \(j\), both evaluated points lie in \(U\cap F\) wherever
\(w>0\), almost surely. We induct on \(k\). With no queries the output
is constant on its domain, so the claim is immediate.

The continuous linear part of the first affine query extends uniquely
from \(F\) to \(H\), and can be written as
\(x\mapsto\langle a,x\rangle\). A zero linear part merely fixes the
first answer and reduces the number of stages. Otherwise put
\[
\begin{aligned}
 \sigma^2&=\langle Ca,a\rangle>0,
 &m(t)&=\frac{Ca}{\sigma^2}t,\\
 H_0&=\ker a,
 &F_0&=F\cap H_0,\\
 C_0&=C-\frac{Ca\otimes Ca}{\sigma^2}.
\end{aligned}
 \tag{W4}
\]
As before, \(Z=m(Y)+Z_0\), with \(Y\) having the normal density \(p\)
and \(Z_0\) an independent centered Gaussian on \(H_0\).
Its covariance \(C_0\) is injective on \(H_0\). Indeed, equality in
the covariance Cauchy--Schwarz inequality would make
\(C^{1/2}h\) proportional to \(C^{1/2}a\); injectivity of \(C\)
then makes \(h\) proportional to \(a\), and \(h\in H_0\) forces
\(h=0\). Moreover \(C_0(H_0)\subset F_0\), and \(Z_0\in F_0\)
almost surely. The subspace \(F_0\) is Borel and dense in \(H_0\):
approximate an element of \(H_0\) by elements of \(F\), then correct
their \(a\)-values using one fixed vector of \(F\) on which \(a\)
does not vanish. These are exactly the hypotheses needed for
induction.

Set
\[
 c_j=\langle a,v_j\rangle\longrightarrow0,\qquad
 v'_j=v_j-m(c_j)\in F_0 .
 \tag{W5}
\]
The \(v'_j\) are weakly null in \(H_0\), and have the same limiting
upper norm as the \(v_j\). After changing \(y=t+c_j\), the shifted
integral in (W3) becomes
\[
 \int_\R\int_{H_0}p(y-c_j)w(m(y-c_j)+z)
       A(m(y)+z+v'_j)\,d\mu_0(z)\,dy .
 \tag{W6}
\]
The function \(f(y,z)=p(y)w(m(y)+z)\) belongs to
\(L^1(dy\,d\mu_0)\). Thus \(f(y-c_j,z)\) can be replaced by
\(f(y,z)\) in (W6), with an error tending to zero, by ordinary
\(L^1\) translation continuity. Boundedness of \(A\) is all that is
used in this replacement.

For each fixed \(y\), let
\[
 U_y=(U-m(y))\cap H_0,\qquad
 A_y(z)=A(m(y)+z),\qquad w_y(z)=w(m(y)+z).
 \tag{W7}
\]
The actual output \(A_y\) is Borel, and the first answer is fixed, so
it is realized by the remaining \(k-1\) affine measurements.
The support of \(w_y\) has distance at least \(R\) from
\(H_0\setminus U_y\), since distance in the affine fiber is at least
ambient distance. Induction therefore gives
\[
 \int_{H_0}w_y(z)\bigl[A(m(y)+z+v'_j)-A(m(y)+z)\bigr]\,d\mu_0(z)
                         \longrightarrow0
 \tag{W8}
\]
for every fixed \(y\). These inner integrals are Borel functions of
\(y\), because they are written using the actual Borel output,
not the possibly nonmeasurable continuation rules. They are bounded
in absolute value by \(2M\norm w_\infty\). Dominated convergence
against \(p(y)\,dy\), together with (W6), proves (W3).
\end{proof}

\subsection{The continuous-information lower bound}

\begin{proof}[Completion of the proof of Theorem~\ref{thm:569}]
The continuous dual of \(F=(\ell_1,\norm{\cdot}_2)\) consists exactly
of restrictions of \(\ell_2\) functionals. Also \(F\) is Borel in
\(H=\ell_2\), since
\[
 \ell_1=\bigcup_{q\in\N}\bigcap_{m\in\N}
       \left\{x\in\ell_2:\sum_{j=1}^m\abs{x_j}\le q\right\}.
 \tag{569.9}
\]
Take \(Z_j=2^{-j}g_j\). Then \(Z\in\ell_1\) almost surely, and its
covariance \(C=\operatorname{diag}(4^{-j})\) is injective and maps
\(\ell_2\) into \(\ell_1\), by Cauchy--Schwarz. Thus (W1) holds.

Fix \(0<r<1\), choose \(0<\rho<1-r\), and condition on
\(\norm Z_2\le\rho\). This event has positive probability: make a
finite Gaussian head small, and make the independent tail small by
first making its expected squared norm small. In
Lemma~\ref{lem:weakshift}, use the corresponding normalized indicator
weight, the open unit ball \(U\) of \(\ell_2\), any
\(r<R<1-\rho\), and \(v_j=re_j\).

First consider a strict finite query budget, and clip every output
to \([0,1]\). This does not increase error, and makes the lemma
applicable. For every fixed seed,
\[
 \E\bigl[A_\omega(Z+re_j)-A_\omega(Z)
                 \mid\norm Z_2\le\rho\bigr]\longrightarrow0.
 \tag{569.10}
\]
Joint measurability of the original output and boundedness allow
averaging (569.10) over an arbitrary seed probability space.
The conditioned Gaussian is symmetric, so
\[
 \E\bigl[S(Z+re_j)-S(Z)\mid\norm Z_2\le\rho\bigr]
       =r^2+2r\,\E[Z_j\mid\norm Z_2\le\rho]=r^2 .
 \tag{569.11}
\]
If \(E_*=\sup_{x\in B}\E_\omega\abs{S(x)-A_\omega(x)}\), the
triangle inequality yields
\[
 r^2\le 2E_*+
       \abs{\E_{Z,\omega}
             [A_\omega(Z+re_j)-A_\omega(Z)\mid\norm Z_2\le\rho]} .
 \tag{569.12}
\]
Let \(j\to\infty\), then \(r\uparrow1\), choosing \(\rho<1-r\)
each time. This proves the expected-absolute-error floor \(E_*\ge1/2\),
and hence the same RMS floor.

For expected cardinality, truncate the clipped algorithm after \(k\)
queries and return zero when it has not stopped. Output and
cardinality measurability make this a jointly Borel strict-budget
algorithm. If the original uniform RMS error is \(e\), then
\[
        \E\abs{S(x)-A^{(k)}_\omega(x)}^2\le e^2+n/k ,
 \tag{569.13}
\]
because \(0\le S(x)<1\). The strict-budget result gives
\(e^2+n/k\ge1/4\) for all \(k\), and hence \(e\ge1/2\).
The constant output \(1/2\) attains that supremal error without
measurements. This proves (569.2).
\end{proof}


\begin{thebibliography}{99}
\bibitem{hnw2011}A. Hinrichs, E. Novak and H. Wo\'zniakowski, \emph{Discontinuous information in the worst case and randomized settings}, Math. Nachr. \textbf{286} (2013), 679--690. Primary manuscript: \url{https://arxiv.org/abs/1106.2945}; \url{https://doi.org/10.1002/mana.201100128}.
\end{thebibliography}
\end{document}
